% ============================================================
\section{Background}
\label{sec:background}
% ============================================================

\subsection{ctDNA Biology and the Liquid Biopsy Signal}

When tumour cells die, they release genomic DNA fragments into the bloodstream.
This circulating tumour DNA (\ctDNA{}) is a subset of the broader pool of
cell-free DNA (cfDNA) shed by all cells in the body.  The physics of cfDNA
fragmentation are informative: because cfDNA is cleaved by nucleosomal wrapping
and endonuclease activity, it is highly fragmented, with a characteristic modal
size of $\sim$167~bp corresponding to mono-nucleosome-protected
DNA~\cite{mouliere2018enhanced}.  Tumour-derived cfDNA fragments tend to be
slightly shorter than non-tumour cfDNA, with a sub-nucleosomal peak at
$\sim$143~bp, providing an additional dimension for tumour signal enrichment.

The fraction of cfDNA that is tumour-derived (\ctDNA{} fraction) varies
enormously across clinical contexts: from $>50\%$ in patients with high
metastatic burden to $<0.01\%$ in early-stage disease or post-treatment MRD
monitoring.  For the most demanding MRD application, only 1 in 10,000 cfDNA
molecules may carry the somatic mutation, meaning that at $3000\times$ sequencing
depth, a typical targeted panel position receives fewer than 1 read from a
mutant molecule.  Achieving reliable detection at this level requires both
ultra-deep sequencing and sophisticated error suppression strategies.

\subsubsection{Fragment Size and K-mer Analysis}

The $\sim$167~bp modal fragment length has a direct impact on k-mer-based
analysis.  For k-mers of length $k = 31$, each cfDNA fragment produces
approximately $167 - 31 + 1 = 137$ overlapping k-mers.  A single-base variant
located in the middle of a fragment contributes up to 31 k-mers that span the
variant position.  At $1000\times$ sequencing depth with $0.5\%$ VAF, the
expected count for each of these 31 k-mers is $1000 \times 0.005 = 5$.  This
is above the minimum count threshold used in the \km{} liquid biopsy
configuration (\texttt{--count 2}), making k-mer-based detection theoretically
feasible at this VAF and depth.

\subsubsection{Clonal Haematopoiesis as a Confounding Factor}

Clonal haematopoiesis of indeterminate potential (CHIP) presents a significant
interpretive challenge in liquid biopsy~\cite{blasco2015chip}.  As individuals
age, haematopoietic stem cells acquire somatic mutations that confer clonal
fitness advantages.  These CHIP clones can reach frequencies of several percent
in the blood, and because cfDNA includes DNA shed by white blood cells, CHIP
mutations can appear as apparent ``tumour'' variants in liquid biopsy assays.
Common CHIP genes (DNMT3A, TET2, ASXL1, TP53) overlap substantially with solid
tumour driver genes, making the distinction between CHIP and true tumour
\ctDNA{} critical.  A panel-of-normals built from age-matched healthy controls
is the most effective strategy for suppressing CHIP-derived false positives in
cfDNA assays~\cite{blasco2015chip}.

\subsection{UMI-Based Error Correction and Duplex Sequencing}

The $\sim$0.1\% per-base sequencing error rate of Illumina platforms is of
the same order of magnitude as the VAF of clinical relevance in MRD monitoring.
Without error correction, it is impossible to distinguish a true variant present
at 0.1\% VAF from a systematic sequencing error at the same rate.

Unique molecular identifiers (UMIs) are short degenerate oligonucleotide barcodes
attached to each DNA molecule during library preparation, before PCR amplification.
Because UMIs are attached to the original template molecule, all PCR copies of
that molecule share the same UMI.  After sequencing, reads sharing the same UMI
and mapping to the same genomic location can be grouped into a \emph{UMI family}
and their sequences compared to produce a consensus that corrects PCR and
sequencing errors~\cite{schmitt2012detection}.

\textbf{Simplex consensus} (single-strand consensus sequencing, SSCS) uses reads
from one strand of each UMI family, achieving error rates of $\sim$10$^{-4}$.
\textbf{Duplex consensus} (DCS) requires agreement from both strands of the same
original DNA molecule (the ``duplex'' formed by the two complementary strands),
achieving error rates of $\sim$10$^{-7}$~\cite{kennedy2014detecting}.  Duplex
sequencing is the gold standard for the lowest-VAF detection but requires higher
sequencing depth (approximately $4\text{--}6\times$ the desired depth of unique
molecule coverage) and increases per-sample cost.

The \kam{} pipeline uses HUMID~\cite{pockrandt2020humid} for UMI-based
deduplication, collapsing UMI families to consensus reads before k-mer counting.
This simplex-level error correction reduces the effective error rate by $10\text{--}100\times$
compared to raw reads, enabling the \texttt{--count 2} threshold to function as
an effective minimum-two-molecule requirement when per-family PCR duplicates are
removed.

\subsection{K-mer Counting: Jellyfish and Alternatives}

\subsubsection{Jellyfish}

\jf{}~\cite{marcais2011fast} is the de facto standard k-mer counter for
bioinformatics applications.  Its key algorithmic innovation is a \emph{lock-free
concurrent hash table} that uses CPU-level compare-and-swap (CAS) instructions
rather than mutex locks to coordinate concurrent k-mer insertion and increment
operations.  This enables near-linear scaling with thread count for the counting
phase, achieving throughput of approximately $10^9$ k-mers per minute on
multi-threaded hardware.

K-mers are stored in 2-bit encoding (A$=00$, C$=01$, G$=10$, T$=11$), so a
31-mer requires only 62 bits and fits in a single 64-bit word.  \jf{} supports
\emph{canonical mode} (\texttt{-C}), in which each k-mer and its reverse
complement are counted together under the lexicographically smaller representation.
This halves memory requirements and automatically handles both strand
orientations, which is essential for double-stranded DNA data.

The \texttt{.jf} file format stores a JSON header (containing k-mer length, hash
table size, canonical flag, and counter configuration) followed by the binary
hash table.  K-mer queries are $O(1)$ through open-address hashing with linear
probing.

\subsubsection{Alternative Counters}

KMC~2~\cite{deorowicz2015kmc2} and DSK~\cite{rizk2013dsk} are disk-based k-mer
counters that scale to extremely large genomes by minimising RAM usage through
external-memory sorting.  For targeted sequencing data with a relatively small
number of distinct k-mers, \jf{}'s in-memory approach is typically faster.
\kmerdet{}'s \texttt{KmerDatabase} trait abstracts the counting backend, allowing
future integration of these alternatives or a native Rust counter without changes
to the variant detection code.

\subsection{K-mer Graphs in Bioinformatics}

The de Bruijn graph formalism~\cite{zerbino2008velvet,bankevich2012spades} provides
the theoretical foundation for k-mer-based sequence analysis.  In a de Bruijn
graph, each node represents a distinct $(k-1)$-mer and each edge represents a
k-mer overlap between adjacent nodes.  Equivalently, nodes can be k-mers
themselves, with edges connecting k-mer pairs where the $(k-1)$-suffix of the
source matches the $(k-1)$-prefix of the target.  De Bruijn graphs are used for
genome assembly (SPAdes~\cite{bankevich2012spades}, Velvet~\cite{zerbino2008velvet}),
error correction (Quake~\cite{kelley2010quake}), and variant calling
(Cortex~\cite{iqbal2012novo}, Scalpel~\cite{narzisi2014accurate}).

The \km{} algorithm builds a \emph{subgraph} of the de Bruijn graph restricted
to k-mers discovered through DFS walking from a target region.  By weighting
reference edges at 0.01 and non-reference edges at 1.0, Dijkstra's shortest-path
algorithm preferentially recovers the reference path, and subsequent removal of
reference edges exposes alternative (variant) paths.  This targeted-graph approach
is computationally cheaper than whole-genome de Bruijn graph assembly while
preserving the ability to detect complex rearrangements such as ITDs.

\subsection{Variant Detection: Alignment-Based vs.\ K-mer-Based}

\subsubsection{Alignment-Based Callers}

Alignment-based somatic variant callers (Mutect2~\cite{cibulskis2013mutect},
Strelka2~\cite{saunders2012strelka}) operate on BAM files produced by alignment
of reads to a reference genome with BWA-MEM~\cite{li2009bwa}.  These tools
construct probabilistic models of read counts, base qualities, mapping qualities,
and strand orientation to compute a posterior probability or likelihood score for
each candidate variant.  Their strengths include rich quality annotations, mature
VCF output, and high sensitivity for SNVs ($>90\%$ at VAF $>0.1\%$).

Their limitations include reference bias (reads that align poorly due to
insertion sequences or structural variants may be filtered as multi-mappers),
sensitivity to MAPQ filtering (particularly problematic in repetitive regions),
and computational cost ($>30$ minutes per sample for targeted panels).

\subsubsection{K-mer-Based Callers}

K-mer-based callers~\cite{pelletier2021kupl,iqbal2012novo,denti2023kmerald,formenti2022merfin}
operate without alignment and are therefore immune to reference bias and
multi-mapping artefacts.  Cortex~\cite{iqbal2012novo} uses coloured de Bruijn
graphs to call variants jointly across multiple samples by differencing k-mer
populations.  The 2-kupl approach~\cite{pelletier2021kupl} detects variants by
identifying k-mers present in a case sample but absent or depleted in a
matched control.  \km{}~\cite{audoux2020km} combines targeted reference walking
with graph pathfinding and NNLS quantification to support both de novo and
tumor-informed detection.

The most significant weakness of k-mer-based callers is reduced INDEL
sensitivity: insertions of length $L \geq k$ cannot produce k-mers spanning
the insertion junction, causing near-complete loss of signal for large INDELs
at k$=31$.  Complex alignment-based callers that perform local graph assembly
(GATK HaplotypeCaller, Scalpel~\cite{narzisi2014accurate}) generally outperform
k-mer walking for INDEL detection, though the gap narrows for smaller indels
in non-repetitive sequence.

\subsection{Statistical Frameworks for Low-VAF Variant Calling}

The statistical challenge of liquid biopsy variant calling can be framed as a
hypothesis test: given observed k-mer counts, does the data provide sufficient
evidence to reject the null hypothesis that all counts arise from sequencing
error?

\subsubsection{Binomial Model}

The simplest framework models the observed variant k-mer count $c$ as drawn
from a Binomial distribution under the null hypothesis of pure sequencing
error~\cite{wei2011snver,koboldt2016statistical}:
\begin{equation}
  P(X \geq c \mid \text{H}_0) = 1 - \mathrm{BinomialCDF}(c - 1,\, D,\, \varepsilon_k)
  \label{eq:binomial}
\end{equation}
where $D$ is the total k-mer coverage at the variant position and $\varepsilon_k$
is the per-k-mer error rate (approximately $k \times \varepsilon_{\text{base}} / 3$
for a specific single-base substitution, with $\varepsilon_{\text{base}} \approx 0.001$
post-UMI deduplication).  The p-value from equation~\eqref{eq:binomial} is
Phred-scaled as $Q = -10 \log_{10} p$ to produce a standard quality score.

\subsubsection{Combining Evidence Across K-mers: Fisher's Method}

For an SNV, up to $k = 31$ k-mer positions span the variant site, each
producing an independent p-value.  Fisher's method~\cite{fisher1925statistical}
combines these:
\begin{equation}
  X = -2 \sum_{i=1}^{m} \ln p_i \;\sim\; \chi^2_{2m}
  \label{eq:fisher}
\end{equation}
where $m$ is the number of k-mer positions contributing signal.  Because
adjacent k-mers share reads, they are not strictly independent; we use an
effective count $m_{\text{eff}} = m \cdot k / l_{\text{read}}$ where
$l_{\text{read}}$ is the typical read length, following the correction of
Jennings et al.~\cite{jennings2017assessing}.

\subsubsection{Negative Binomial Overdispersion}

Count data from sequencing consistently exhibit overdispersion relative to the
Poisson model, due to PCR amplification bias, GC-dependent capture efficiency,
and coverage non-uniformity~\cite{love2014deseq2}.  The negative binomial model
extends the Poisson by allowing variance $= \mu + \alpha\mu^2$ where $\alpha$ is
the overdispersion parameter.  In practice, $\alpha$ is estimated from the
reference k-mer count distribution along the target region, which should be
approximately uniform under the null hypothesis of no variant.

\subsubsection{Bootstrap Confidence Intervals on rVAF}

The NNLS decomposition produces a point estimate of rVAF but no uncertainty
bound.  A parametric bootstrap resamples k-mer counts from the estimated
distribution (Poisson with mean equal to observed count) and re-solves the
NNLS for each of $B = 1000$ resamples.  The 2.5th and 97.5th percentiles of the
resulting rVAF distribution form the 95\% confidence interval, providing a
clinically interpretable range for longitudinal VAF tracking~\cite{jennings2017assessing}.
